\section{October 7, 2026}
% Based on pages 10--13 of 6-7-week.pdf.

Consistency measures how well exact samples satisfy a multistep
formula. To obtain convergence, we must also control how the
recurrence propagates perturbations. This leads to zero-stability
and the root condition. We continue to use a uniform grid,
\(t_n=t_0+n\tau\), and the polynomials \(\rho\) and \(\sigma\)
from \eqref{eq:lmm-polynomials-oct-05}.

\subsection{A consistent method that is unstable}

Consider the explicit two-step formula from class,
\begin{equation}\label{eq:unstable-lmm-oct-07}
 x_{n+2}+4x_{n+1}-5x_n=\tau(4f_{n+1}+2f_n).
\end{equation}
Its characteristic polynomials are
\[
 \rho(\zeta)=\zeta^2+4\zeta-5=(\zeta-1)(\zeta+5),
 \qquad \sigma(\zeta)=4\zeta+2.
\]
Since \(\rho(1)=0\) and \(\rho'(1)=6=\sigma(1)\), the method
is consistent. In fact, the order conditions for \(m=2,3\) hold:
\begin{align*}
 4+2^2&=2(4),\\
 4+2^3&=3(4).
\end{align*}
At \(m=4\), however, \(4+2^4=20\neq16=4(4)\).
Thus its order is exactly three, with residual
\[
 \mathcal L_\tau[x](t_n)
       =\frac16\tau^4x^{(4)}(t_n)+O(\tau^5).
\]

Now apply it to \(x'=0\). The exact solution is constant, but the
numerical recurrence is
\[
 x_{n+2}+4x_{n+1}-5x_n=0.
\]
Trying \(x_n=\zeta^n\) gives \(\rho(\zeta)=0\), so every
solution is of the form
\begin{equation}\label{eq:unstable-mode-oct-07}
 x_n=a+b(-5)^n.
\end{equation}
The constant mode is correct; the second mode magnifies perturbations
exponentially.

For a concrete failure of convergence, take the IVP \(x'=0\),
\(x(0)=0\), and starting values \(x_0=0\), \(x_1=\tau^3\).
These starting errors tend to zero at third order, but
\[
 x_n=\frac{\tau^3}{6}\bigl(1-(-5)^n\bigr).
\]
At fixed \(T>0\), with \(n=T/\tau\), its magnitude grows without
bound along the grids \(\tau=T/N\). Small consistency error alone
does not ensure convergence.

\subsection{The root condition and zero-stability}

Setting \(f=0\) in an LMM leaves
\begin{equation}\label{eq:zero-recurrence-oct-07}
 \rho(E)x_n=0.
\end{equation}
A simple nonzero root \(\zeta\) of \(\rho\) produces a mode
\(\zeta^n\); a root of multiplicity \(m\) produces modes
\[
 \zeta^n,\quad n\zeta^n,\quad\ldots,\quad n^{m-1}\zeta^n.
\]
Roots strictly inside the unit disk give bounded, decaying modes,
even when repeated. A zero root corresponds to a finite transient
in the companion-matrix description. On the unit circle, any
polynomial factor in \(n\) destroys boundedness.

\begin{definition}[Root condition and zero-stability]
 An LMM is \emph{zero-stable} if every root of \(\rho\) satisfies
 \begin{equation}\label{eq:root-condition-oct-07}
  |\zeta|\leq1,
  \qquad\text{and every root with }|\zeta|=1\text{ is simple}.
 \end{equation}
 This is the \emph{root condition}.
\end{definition}

Equivalently, there is a bound independent of the number of steps,
\[
 \sup_{n\geq0}\|x_n\|
       \leq K\max_{0\leq j<k}\|x_j\|,
\]
for solutions of the homogeneous recurrence
\eqref{eq:zero-recurrence-oct-07}. This uniformity matters because
the number of steps to a fixed time tends to infinity as
\(\tau\to0\). The term ``zero'' refers to the test equation
\(x'=0\); absolute stability at nonzero \(z=\tau\lambda\) will
require a different polynomial.

\begin{example}[A repeated unit root]
 If \(\rho(\zeta)=(\zeta-1)^2\), then
 \[
  x_{n+2}-2x_{n+1}+x_n=0,
  \qquad x_n=a+bn.
 \]
 Thus \(|\zeta|\leq1\) by itself is not enough. Starting values
 \(x_0=0\), \(x_1=\tau\) yield \(x_N=T\) at \(N\tau=T\),
 despite both starting values tending to zero.
\end{example}

\begin{example}[The methods from October 5]
 Leapfrog and Simpson's multistep formula have
 \(\rho(\zeta)=\zeta^2-1\). Both roots, \(1\) and \(-1\),
 are simple, so both methods are zero-stable. The \(-1\) mode
 alternates without decay; zero-stability permits such a bounded
 parasitic mode.
\end{example}

\subsection{Convergence and the Dahlquist equivalence theorem}

\begin{definition}[Convergence of an LMM]
 An LMM is convergent if, for every sufficiently regular IVP and
 every fixed finite interval \([t_0,T]\), starting values with
 \[
  e_{\rm start}(\tau)
   :=\max_{0\leq j<k}\|x_j-x(t_j)\|\longrightarrow0
 \]
 give numerical solutions satisfying
 \[
  \max_{0\leq n\leq N}\|x_n-x(t_n)\|\longrightarrow0
  \quad\text{as }\tau\to0,\qquad t_N\leq T.
 \]
\end{definition}

Because the finitely many starting times approach \(t_0\), the
starting condition is equivalently \(x_j\to x(t_0)\) for
\(j=0,\ldots,k-1\). Keeping \(x(t_j)\) in the definition makes
the required accuracy easier to quantify.

\begin{theorem}[Dahlquist]
 A linear multistep method is convergent if and only if it is
 consistent and zero-stable, under the usual smoothness and
 Lipschitz assumptions on the IVP.
\end{theorem}

More quantitatively, for a zero-stable method of consistency order
\(p\), sufficiently small \(\tau\) gives an estimate of the form
\begin{equation}\label{eq:lmm-global-bound-oct-07}
 \boxed{\max_{0\leq n\leq N}\|x_n-x(t_n)\|
  \leq C e^{L_*(T-t_0)}
          \bigl(e_{\rm start}(\tau)+\tau^p\bigr).}
\end{equation}
The constants are independent of \(\tau\), but may depend on
the method, the interval, and bounds on the solution and \(f\).
In particular, order-\(p\) convergence requires starting errors
of size \(O(\tau^p)\).

To see the mechanism, put \(e_n=x_n-x(t_n)\). Subtracting the
exact residual equation from the numerical formula gives
\begin{equation}\label{eq:lmm-error-recursion-oct-07}
 \sum_{j=0}^k\alpha_je_{n+j}
 =\tau\sum_{j=0}^k\beta_j
       \bigl(f(t_{n+j},x_{n+j})-f(t_{n+j},x(t_{n+j}))\bigr)
       -\mathcal L_\tau[x](t_n).
\end{equation}
The root condition bounds propagation by the homogeneous recurrence.
The Lipschitz condition controls the first term on the right, and
a discrete Gronwall argument gives a bound in terms of
\[
 e_{\rm start}(\tau)
       +\sum_{n=0}^{N-k}\|\mathcal L_\tau[x](t_n)\|.
\]
Since each residual is \(O(\tau^{p+1})\) and there are
\(O(\tau^{-1})\) steps, their sum is \(O(\tau^p)\).
This is the multistep counterpart of the local-to-global argument
for one-step methods on August 31.

For an implicit formula, the new value solves
\[
 x_{n+k}=\text{known terms}
              +\tau\frac{\beta_k}{\alpha_k}f(t_{n+k},x_{n+k}).
\]
For example, a Lipschitz constant \(L\) gives a contraction when
\[
 \tau\left|\frac{\beta_k}{\alpha_k}\right|L<1.
\]
This supplies a sufficient small-step solvability condition in the
convergence argument.

\subsection{Adams methods}

The Adams family uses only the difference between the two newest
solution values on the left:
\begin{equation}\label{eq:adams-form-oct-07}
 x_{n+k}-x_{n+k-1}
       =\tau\sum_{j=0}^k\beta_jf_{n+j}.
\end{equation}
Thus
\begin{equation}\label{eq:adams-rho-oct-07}
 \rho(\zeta)=\zeta^k-\zeta^{k-1}
            =\zeta^{k-1}(\zeta-1).
\end{equation}
The root at one is simple; all remaining roots are zero. Every
method with this left-hand side is therefore zero-stable, independent
of the derivative weights. Consistency additionally requires
\[
 \sigma(1)=\rho'(1)=1,
 \qquad\text{that is,}\qquad \sum_{j=0}^k\beta_j=1.
\]
The usual Adams weights are obtained by integrating an interpolant
of the derivative values. Excluding the new endpoint derivative
gives Adams--Bashforth; including it gives Adams--Moulton. Their
construction and their nonzero-\(z\) stability are the topics of
the next lecture.
